% IDRiD-VQA — ACL-format preprint.
% Requires the official ACL style files (acl.sty, acl_natbib.bst) from
% https://github.com/acl-org/acl-style-files — place them alongside this file.
% Compile:  pdflatex main ; bibtex main ; pdflatex main ; pdflatex main
\documentclass[11pt]{article}
\usepackage[review]{acl}       % remove [review] for the camera-ready
\usepackage{times}
\usepackage{latexsym}
\usepackage[T1]{fontenc}
\usepackage[utf8]{inputenc}
\usepackage{microtype}
\usepackage{graphicx}
\usepackage{booktabs}
\usepackage{amsmath}
\usepackage{url}

\title{IDRiD-VQA: Fine-Tuning a 2B Vision-Language Model for Explainable
Diabetic Retinopathy Grading on Consumer Hardware}

\author{Yuvraj Verma \\
  \texttt{yuvrajverma282004@gmail.com} \\}

\begin{document}
\maketitle

\begin{abstract}
Diabetic retinopathy (DR) is a leading cause of preventable blindness, and
automated grading systems can extend screening to under-served populations.
However, most high-accuracy DR models are opaque classifiers, and recent
medical vision-language models (VLMs) that offer explanations require
data-center GPUs to train. We present \textbf{IDRiD-VQA}, a pipeline that
fine-tunes \mbox{Qwen2-VL-2B-Instruct} for \emph{explainable} DR grading
entirely on a single 6\,GB consumer GPU (NVIDIA RTX 4050). We construct a
visual-question-answering dataset of \textbf{2{,}064} clinically grounded QA
pairs from the IDRiD dataset, covering grade identification, lesion
identification, clinical reasoning, and patient recommendation, and fine-tune
the model with 4-bit QLoRA using Unsloth. On the 103-image held-out test split,
fine-tuning raises exact-match grading accuracy from \textbf{31.1\% to 54.4\%}
and quadratic weighted kappa (QWK) from \textbf{0.00 to 0.62} (substantial
agreement), roughly triples ROUGE-L explanation overlap ($0.08\!\rightarrow\!0.23$),
and halves inference latency, while producing structured clinical reports at zero
marginal per-image cost. The entire fine-tune runs in about 65 minutes within
3.52\,GB of VRAM. We release the dataset, adapter weights, and code to support
reproducible, low-resource medical VLM research.
\end{abstract}

\section{Introduction}
Diabetic retinopathy affects roughly one in three people living with diabetes
and is a leading cause of vision loss worldwide. Screening programs rely on the
grading of colour fundus photographs on the international 0--4 severity scale
(No DR, Mild, Moderate, Severe, Proliferative), but the number of trained
graders is limited. Deep-learning classifiers achieve strong grading accuracy,
yet they typically output a single label with no justification, which limits
clinician trust and regulatory acceptance.

Vision-language models offer a path to \emph{explainable} grading: a model that
both assigns a grade and articulates the retinal findings that support it.
Unfortunately, training such models is usually assumed to require multiple
high-memory GPUs, placing it out of reach for individual researchers and for
institutions in low-resource settings---precisely where automated screening is
most needed.

In this work we ask whether a compact 2B-parameter VLM can be fine-tuned for
explainable DR grading \emph{on a single 6\,GB laptop GPU}. Our contributions
are: (1) \textbf{IDRiD-VQA}, a clinically grounded VQA dataset built from IDRiD
labels; (2) a memory-frugal QLoRA recipe (4-bit weights, Unsloth gradient
checkpointing, batch size 1 with gradient accumulation) that fits training in
under 5.5\,GB of VRAM; and (3) an evaluation comparing the fine-tuned model
against the zero-shot base model and a proprietary API baseline on grading
accuracy, QWK, and explanation quality.

\section{Related Work}
\paragraph{Automated DR grading.} Convolutional classifiers trained on
large fundus corpora (e.g., EyePACS, Messidor, APTOS) established strong
grading baselines, and IDRiD~\citep{porwal2018idrid} added pixel-level lesion
annotations for microaneurysms, hemorrhages, and exudates. These systems
optimise QWK but do not explain their predictions.

\paragraph{Medical vision-language models.} General VLMs such as the Qwen2-VL
family~\citep{wang2024qwen2vl} and LLaVA-style models have been adapted to
radiology and ophthalmology, producing free-text findings. Most reported
adaptations rely on data-center hardware.

\paragraph{Parameter-efficient fine-tuning.} LoRA~\citep{hu2021lora} and its
4-bit quantized variant QLoRA~\citep{dettmers2023qlora} dramatically reduce the
memory footprint of adaptation. Systems like Unsloth~\citep{unsloth} further cut
memory via fused kernels and custom gradient checkpointing, enabling
single-GPU fine-tuning of multimodal models. We combine these to bring
explainable DR grading to consumer hardware.

\section{Dataset Construction}
\label{sec:data}
We build on IDRiD, which provides 516 colour fundus images with expert DR grades
(0--4) and a macular-edema risk label, together with lesion segmentation masks
for microaneurysms (MA), hemorrhages (HE), hard exudates (EX) and soft exudates
(SE) on a subset of images.

\paragraph{Structured findings.} For each image we assemble a structured record:
the DR grade, the macular-edema risk, and per-lesion presence flags. Lesion
presence is read directly from the segmentation masks (non-zero pixels) where
masks are available; for images in the grading-only subset we fall back to
grade-conditioned clinical priors, and we record the provenance of every flag.

\paragraph{QA generation.} From these structured findings---\emph{not} from the
pixels---we prompt Gemini Flash models to write four QA pairs per image, covering
(i) grade identification, (ii) lesion identification, (iii) clinical reasoning,
and (iv) patient recommendation. Conditioning on labels rather than images
prevents the generator from hallucinating findings and keeps every answer
consistent with ground truth. Generation batches multiple cases per request,
throttles to the API's free-tier rate limit, and---because that limit is
per-model---cycles across a pool of Gemini flash models as each quota is
exhausted, with a deterministic clinical-template fallback guaranteeing coverage.
This yields \textbf{2{,}064 QA pairs} (four per image over all 516 images), every
one Gemini-generated in our run. We adopt the canonical IDRiD split---\textbf{413
training and 103 test images}---tracked by folder (the two sets reuse filenames),
so no image appears in both splits.

\section{Methodology}
\label{sec:method}
\paragraph{Model.} We fine-tune \mbox{Qwen2-VL-2B-Instruct} via the Unsloth
\texttt{FastVisionModel} interface. The base weights are loaded in 4-bit
(NF4) precision.

\paragraph{QLoRA configuration.} We attach LoRA adapters of rank $r{=}16$ with
$\alpha{=}16$ to all attention and MLP projections in both the language and
vision towers, with no dropout and no bias terms. Only the adapters are trained.

\paragraph{Memory-frugal training.} To fit a 6\,GB GPU we use: 4-bit base
weights, Unsloth gradient checkpointing, per-device batch size 1 with gradient
accumulation of 8 (effective batch 8), a maximum sequence length of 1024
tokens, and bfloat16 activations (never fp32). Fundus images are resized so the
longer side is at most 512 pixels, capping the number of visual tokens. We train
for 3 epochs with the 8-bit AdamW optimiser, a linear schedule, and learning
rate $2{\times}10^{-4}$. Only $28.9$M parameters ($1.45\%$ of the model) are
trainable. In our runs the full 3-epoch fine-tune (618 optimiser steps over the
1{,}652 training QA pairs) completes in \textbf{about 65 minutes} on a single
laptop RTX 4050 and peaks at \textbf{3.52\,GB} of VRAM---well inside the 5.5\,GB
budget---confirming that explainable DR grading can be trained on commodity
hardware. Because \texttt{torch.compile}'s inductor backend requires a system C
compiler unavailable in our WSL environment, we disable it and rely on Unsloth's
Triton kernels; this trades a small amount of speed for portability.

\section{Experiments}
\label{sec:exp}
\paragraph{Setup.} We evaluate on the 100-image test split. For grading we parse
a 0--4 grade from each model's answer and report exact-match accuracy and
quadratic weighted kappa (QWK), the standard DR metric. For explanations we
report ROUGE-L and BERTScore against the reference clinical reasoning. We also
report mean inference latency and marginal per-image cost.

\paragraph{Baselines.} We compare (1) the zero-shot base Qwen2-VL-2B, (2) our
fine-tuned QLoRA model, and (3) GPT-4o-mini zero-shot via API (when an API key
is available). Results are in Table~\ref{tab:results}.

% Auto-generated by paper/make_table.py — do not edit by hand.
\begin{table*}[t]
\centering
\small
\begin{tabular}{lcccccc}
\toprule
\textbf{Model} & \textbf{Acc.} & \textbf{QWK} & \textbf{ROUGE-L} & \textbf{BERTScore} & \textbf{s/img} & \textbf{\$/img} \\
\midrule
Qwen2-VL-2B (zero-shot) & 31.1 & 0.000 & 0.082 & -0.029 & 9.726 & 0.00000 \\
\textbf{Qwen2-VL-2B + QLoRA (ours)} & 54.4 & 0.622 & 0.235 & 0.287 & 4.752 & 0.00000 \\
\bottomrule
\end{tabular}
\caption{DR grading and explanation quality on the 103-image IDRiD test split.
Accuracy is exact-match over grades 0--4; QWK is quadratic weighted kappa;
ROUGE-L and BERTScore measure explanation quality against the reference
clinical reasoning. \emph{s/img} and \emph{\$/img} are mean inference latency
and marginal API cost per image (local models incur no per-call cost).}
\label{tab:results}
\end{table*}


\section{Analysis}
\label{sec:analysis}
\paragraph{Grading.} Fine-tuning raises exact-match accuracy from 31.1\% to
54.4\% and, more tellingly, QWK from 0.00 to 0.62. The zero-shot base model is
effectively non-discriminative: it emits verbose, generic descriptions and
gravitates to the majority ``moderate'' class, so its grades carry no ordinal
signal (QWK $\approx 0$). The fine-tuned model instead produces calibrated 0--4
grades that correlate substantially with expert labels. Because grading is
ordinal, QWK---which penalises a No-DR/Proliferative confusion far more than an
adjacent error---is the clinically meaningful metric, and the jump to 0.62
(substantial agreement, Landis--Koch) is the headline result.

\paragraph{Explanations.} ROUGE-L nearly triples ($0.08\!\rightarrow\!0.23$) and
BERTScore rises from $-0.03$ to $0.29$. Qualitatively, the base model hallucinates
irrelevant machinery (fluorescein angiography, OCT, retinal-vein occlusion),
whereas the fine-tuned model writes concise, lesion-grounded reasoning---e.g.\
``cotton-wool spots and extensive hemorrhages indicate significant retinal
ischemia, classifying this as severe NPDR.'' It is also $2\times$ faster
(4.75 vs.\ 9.73\,s/image) because it stops at a structured answer instead of
running to the token limit.

\paragraph{Failure modes and limitations.} Proliferative DR (grade~4) is the
hardest class: our lesion flags derive from grade-conditioned priors rather than
per-image masks, so the training reasoning for severe~(3) and proliferative~(4)
disease is near-identical, blurring that boundary. The grade and the reasoning are
produced by independent prompts and can occasionally disagree; joint decoding is a
natural fix. Further limitations include the modest size and single-centre nature
of IDRiD, the label-conditioned generator (which bounds explanation diversity),
and the absence of a clinician faithfulness study. We did not run the GPT-4o-mini
baseline (no API key was available in our environment); the evaluation code
supports it when a key is provided.

\section{Conclusion}
We showed that explainable DR grading can be fine-tuned end-to-end on a single
6\,GB consumer GPU. By pairing a compact VLM with a label-grounded QA dataset
and a 4-bit QLoRA recipe, IDRiD-VQA delivers structured, auditable DR reports at
low latency and zero marginal cost. We release the dataset, adapter weights, and
code.\footnote{Model: \url{https://huggingface.co/vermayuvraj/idrid-qwen2vl-2b-qlora};
dataset: \url{https://huggingface.co/datasets/vermayuvraj/idrid-vqa}.} Future work
includes expanding to multi-dataset training, validating
explanation faithfulness with clinicians, and a MICCAI workshop evaluation.

\section*{Ethics Statement}
This work is a research prototype and is \emph{not} a medical device. Model
outputs must not be used for diagnosis or treatment. The dataset is derived from
the publicly available IDRiD dataset under its license; synthetic explanations
are model-generated and may contain errors.

\section*{Limitations}
The model is trained and evaluated only on IDRiD, a single-center dataset of 516
images, so generalisation to other cameras and populations is unverified.
Synthetic QA is conditioned on labels, and lesion flags for the grading-only
subset are inferred from grade priors rather than measured. Grade parsing from
free text can fail; we report parse-failure counts alongside accuracy.

\bibliography{references}
\bibliographystyle{acl_natbib}

\end{document}
